% page 178 du fac-simile (image p-177 du scan)
% bloc 162.6 x 242.0 mm ; marge gauche 8.0 mm
\pgc{-12.700mm}{0.000mm}{
\l{\cel{3}168}
\l{}
\l{}
\l{\sou{flanko}. I. (\sou{anat.}) Parto latera di la ventro, situita, an}
\l{\cel{3}singla latero di la korpo, inter la lakuno kostala e la}
\l{\cel{3}loko ube komencas la hancho. - II. (\sou{fortifikajo)}. Parto}
\l{\cel{3}di la bastiono, inter la kurteno e la shutro-angulo.-DEFIR}
\l{}
\l{\sou{flarar}. (\sou{trans.}) Probar dicerno segun la odoro. (\sou{anke meta}-}
\l{\cel{3}\sou{fore)}. F.}
\l{}
\l{\sou{flatar.}\cel{1}- (\sou{trans}.)Probar sedukto per laudi, nemeritita od exa-}
\l{\cel{3}jerita. - DEF.}
\l{}
\l{\sou{flatuar}. (\sou{netrans}.)\cel{2}Evakuar gaso intestinala, tra la sfink-}
\l{\cel{3}tero anusala. - DEFIS.}
\l{}
\l{\sou{flatulenta}. (\sou{medic}.) Che qua gasi, venti, akumulesis en la}
\l{\cel{3}kanalo digestala.}
\l{}
\l{\sou{flava}. Di qua la koloro esas olta di oro, di safrano, di ci-}
\l{\cel{3}trono, e c. - FL.}
\l{}
\l{\sou{flebito}. (\sou{patol}.) Inflamuro di la membrano interna di la veini.}
\l{\cel{3}- DEFIS.}
\l{}
\l{\sou{flecho.}\cel{2}I. Armo, ek stango ligna, garnisita, an un ek la ex-}
\l{\cel{3}tremaji, ye ferpeco, pintoza, an la altra ye aleti ek plumi}
\l{\cel{3}o metalo, e quan onu lansas per arko od arbalesto. - II.}
\l{\cel{3}(\sou{arkitekt.}) Alteso vertikala, de la punto maxim alta di vul-}
\l{\cel{3}to, di arko, super la komenco-punti. - III. (\sou{fortifikajo}).}
\l{\cel{3}Masivo salianta angule. - IV. (\sou{geom.}) Perpendiklo de la mezo}
\l{\cel{3}di arko anla kordo qua klozas lu ; o disto di la kordo di}
\l{\cel{3}arko an la tangento paralela. - FIS.}
\l{}
\l{\sou{flegar}. (\sou{trans.}) I. Sorgar atencoze por la komforto di malado,}
\l{\cel{3}di infanteto o di infanto. - II. Striglagar (kavalo).- D.}
\l{}
\l{\sou{flegmo}.\cel{2}Karaktero di la persono kalma, qua su dominacas.-DEFIRS}
\l{}
\l{\sou{flegmono.}\cel{1}(\sou{patol.}) Centropusifanta qua rezultas de la inflameso}
\l{\cel{3}di la tisuo celuloza subpela. - DEFIRS.}
\l{}
\l{\sou{flereto}. Speco di espado, di qua la lamo esas quadrato-seciona,}
\l{\cel{3}tenua e flexebla, generale kun buleto ye la pinto, garnisita}
\l{\cel{3}ye felo, uzata por exercar su pri skermo. - DFIS.}
\l{}
\l{\sou{flexar}. (\sou{trans.}) I. Kurvigar per preso qua interproximigas\cel{2}la}
\l{\cel{3}du extremaji. - II. (\sou{metaf}.)Cesar pokope sua rezisto. -EFIS.}
\l{}
\l{\sou{flexiono.}\cel{1}(\sou{gram.}) Modifikeso di vorto, da la chanjo di lua de-}
\l{\cel{3}zinenci (en la deklino, la konjugo, che ula lingui nacio-}
\l{\cel{3}nala). - DEFIS.}
\l{}
\l{\cel{1}\sou{flintglaso}.\cel{2}Kristalo pura, por la lensi akromata, kompozita}
\l{\cel{3}ek plumbo-silikato e kalio-silikato, e tre refraktiva. DEFRS}
}
